\section{Methodology}

\subsection{Image Retrieval and Dataset Construction}

Re-LAION-2B-en-research-safe was processed across four uniformly sampled shards,
yielding an initial pool of 65,552,809 URL-caption pairs, of which 43,491,938 images
were successfully retrieved after a link decay rate of 33.7\%. CLIP ViT-L/14 embeddings
were computed for all valid images, L2-normalised, and indexed per shard using FAISS
IndexFlatIP for exact cosine similarity retrieval. A set of 200 occupations aligned with
the ISCO-08 taxonomy was constructed to ensure systematic coverage across skill levels
and occupational domains. For each occupation, the top candidates were retrieved using
the template \texttt{"A photo of a/an \{profession\}"} with gender-neutral phrasing,
filtered for semantic relevance via a negative prompt dominance criterion, and passed
through a multi-stage structural validation pipeline requiring exactly one detected face
and person, verified spatial association between them, and successful MediaPipe FaceMesh
landmark extraction. YOLOv12l-Face was selected as the face detector following
evaluation on WIDER Face, achieving the highest AP@0.5 of 0.905 and the best small-face
recall of the four evaluated models. After structural filtering, the top 1,000
highest-scoring images per occupation were retained, yielding a balanced
profession-conditioned subset of 200,000 images with a mean structural acceptance rate
of 33.51\%.

A corresponding 200,000-image generation corpus was produced using Stable Diffusion
v1.5 with LCM-LoRA acceleration, selected from three LAION-traceable model
configurations on throughput grounds (1.5 s/image versus 5.6 s for standard SDv1.5).
The generation prompt \texttt{"full-body realistic photo of a \{profession\}, standing, centred, in a professional setting appropriate for their occupation, sharp focus, 35mm lens, natural professional lighting"} was used throughout, with classifier-free guidance fixed at $s=1.0$ to establish a conservative lower bound on demographic bias free from
inference-time amplification. Generated images passed through the same structural
validation pipeline as retrieved images, achieving a mean acceptance rate of 66.38\%.

\subsection{Demographic Annotation}

All images underwent face segmentation via MediaPipe FaceMesh, producing tightly
cropped, background-removed face crops to prevent background shortcut learning. Gender
was annotated using the Realistic Gender Classifier, selected from five candidate models
for its highest overall accuracy on CCv2 face-only crops (88.6\%) and the smallest
male-female accuracy gap (6.2 pp) relative to alternatives exhibiting gaps of up to
40.2 pp. Skin tone was estimated using a VGG-16 MST10 regression model trained on CCv2,
selected via a factorial ablation across output formulation, colour space, label
granularity, and background treatment. Predictions were collapsed into Light (MST~1--3),
Mid (MST~4--7), and Dark (MST~8--10) bins for aggregate analysis. All demographic
labels are treated as model-derived proxies of perceived attributes rather than ground
truth. Observed distributions were contextualised against BLS occupational statistics
as the primary real-world reference, supplemented by ILOSTAT global data. A
race-to-skin-tone mapping via FairFace was derived but subsequently abandoned due to
mid-bin collapse in the conditional distributions, limiting its utility for real-world
comparison.

\subsection{Image Transformation Pipeline}

The transformation pipeline serves two probing objectives: dataset classification,
following Zeng et al., which tests whether a ConvNeXt-Tiny classifier can distinguish
Re-LAION-5B, SDv1.5, and COCO images from each transformation alone; and demographic
artifact probing, following Meister et al., which tests whether gender and skin tone can
be predicted from isolated visual channels. The three-class design, extending Zeng et
al.'s two-dataset setup by adding COCO as a structurally independent reference, enables
per-class accuracy attribution from the confusion matrix, localising which dataset
carries each type of spurious signal. For demographic probing, simple probe models are
used deliberately: logistic regression for low-dimensional representations such as mean
RGB and object detection features, and ResNet-50 for spatial inputs, following Meister
et al.

The transformation hierarchy spans three levels. At the structural level, monocular
depth maps were produced using Depth-Anything-V2 ViT-S; Canny edge maps were derived
using classical edge detection with hysteresis thresholds matching Zeng et al.; and
frequency domain representations were produced as low-pass and high-pass filtered
images via a hard circular FFT mask at radius $r=40$. Pose was estimated using
YOLOv8l-Pose, with COCO-17 keypoints normalised for scale and position into a
34-dimensional vector as direct probe input. At the spatial-statistical level, pixel
shuffle, patch shuffle at four scales (2$\times$2 to 16$\times$16), and mean RGB
transformations progressively isolate global colour statistics from local and global
spatial structure. At the semantic level, segmentation maps were produced using
Mask2Former with a Swin-Large backbone; object detection representations were produced
using YOLOv8x with OpenImages V7 weights, with demographic ambiguity-inducing categories
remapped to neutral equivalents; image captions were generated using BLIP Large,
selected via multi-metric evaluation on COCO and Flickr30k, with gender-neutral
remapping applied post-generation; and VAE reconstructions were produced using the
Stable Diffusion v1.5 KL autoencoder.

Occlusion variants followed Meister et al.'s protocol, with five conditions per image:
full image with no background (Full NoBg), person mask replaced by white fill
(MaskSegm), person silhouette on black background (MaskSegm NoBg), rectangular bounding
box replaced by white fill (MaskRect), and rectangular bounding box on black background
(MaskRect NoBg). Person masks were predicted by Mask2Former COCO, selected for its
highest detection rate of 96.5\% among five evaluated segmentation models.

\subsection{Inference-Time Debiasing Pipeline}

A 10,000-image debiased corpus was generated for the doctor profession, concentrating
images on a single high-skew occupation to achieve sufficient statistical power for
detecting demographic shifts. The pipeline composes three inference-time interventions
applied to a frozen SDv1.5 with classifier-free guidance at $s=7.0$, deliberately
activating the model's full demographic prior to provide a demanding test of the
debiasing methods.

ITI-GEN~\cite{Zhang_2023_ITIGEN} inclusive token embeddings were trained on CCv2 across
gender (female, male) and skin tone (MST~1--10), producing 20 demographic combinations
each generating 500 images. Tokens were applied via the \texttt{prompt\_prepend}
mechanism, concatenating inclusive embeddings with the CLIP text encoding. ControlNet
OpenPose~\cite{LZhang_2023_ControlNet} was applied at a conditioning scale of 0.6 using
a single fixed canonical pose skeleton, enforcing structural consistency across generated
images and preventing pose-correlated demographic variation from confounding the
annotation pipeline. Chamfer-based colour alignment~\cite{Shum_2025_ColorAlignment} was
applied to the predicted clean sample $\hat{x}_0$ at denoising timesteps $t > 600$,
steering the colour distribution toward a skin-tone-neutral mean RGB reference derived
from the standard SDv1.5 doctor corpus. LCM-LoRA was excluded because its reduced step
count combined with per-step Chamfer modifications produced insufficiently coherent
outputs. SW-Guidance was evaluated but excluded on throughput grounds, as repeated
decode-gradient-encode operations at each step rendered it impractical at the required
scale.